\section{Risk Management}


\begin{frame}{Quote of the Day}
    
    \mytext{
    \begin{quote}
        La victoire va à celui qui prend le plus de risques.
        
        \flushright Jean-Claude \textsc{Killy} (ski racer)
    \end{quote}
    }
\end{frame}


\begin{frame}{Risk Management}
    
    \mytext{
    \begin{itemize}
        \item \textbf{What could possibly go wrong?}
        \item Risk management $\neq$ issue managment
        \item Foresight of alternative futures
        \item Planning alternative futures
    \end{itemize}
    }
\end{frame}

\begin{frame}{Risk Definition}
    
    \mytext{
    \begin{itemize}
        \item ISO 31000
        \item \textbf{effect of uncertainty on objectives}
        \item \emph{not just chance or probability of loss}
        \item Stakeholders: \emph{Person or persons that can affect, be affected by, or perceive themselves to be affected by a decision or activity}
        \item Formalise risk ownership
    \end{itemize}
    }
\end{frame}


\begin{frame}{Risk and Payoff -- "No Risk, no fun!"}
    
    \mytext{
    
    Economic theory:

    \begin{itemize}
        \item risk must be \textbf{compensated}
        \item[] ... for the same expected outcome $E$
        \item people are risk-averse
        \item risk must be \emph{over}compensated
        \item risk correlated to returns
    \end{itemize}
    
    Note: risk $\neq$ uncertainty!
    
    }
\end{frame}




\begin{frame}{Sources of Risk}
    
    \mytext{
    \begin{itemize}
        \item Every process and result and person bears risk
        \item Interviews and questionnaires for research
        \item There are multiple types of risk
        \item by origin / area
        \item by variance / whether known or not
        \item Plan for eventualities
        \item Contingency plan for \emph{exeptional risk}
    \end{itemize}
    }
\end{frame}


\begin{frame}{Types of Risk by Origin / Area}
    
    \mytext{
    \begin{itemize}
        \item Human resources risk
        \item Inventory risk
        \item Information risk
        \item Data risk
        \item Technological risk
        \item Financial: sources
        \item Financial: cash flow
        \item Financial: costs\end{itemize}
    }
\end{frame}

\begin{frame}{Types of Risk by Origin / Area}
    
    \mytext{
\begin{itemize}
        \item Market risks
        \item Planning risk
        \item Legal risk: contracts
        \item Regulatory risk / political risks
        \item Personal risks
        \item \textbf{Stakeholder risks}
        \item Risks as opportunities
        \item Commitment and confidence
        \item[] (leadership risk)
    \end{itemize}
    }
\end{frame}

\begin{frame}{Commitment: Game Theory: Stag Hunt}
    \showmytikz{../shared/tikz/tikz-PM-stag-hunt}
    \mycomment{
    

    \begin{itemize}
        \item both go hunting: best for both
        \item individually: sure result if gather
        \item each has an incentive to chicken out / let the other down
    \end{itemize}

    How to make people commit?
    }
\end{frame}

\begin{frame}{Risk Measurement}
    
    \mytext{
    \begin{itemize}
        \item Are there \emph{alternative outcomes}
        \item ...or is there a probability distribution
        \item Standard derivation $\sigma$
        \item Known risks are \textbf{insurable} risks
        \item Non-insurable risk is called \textit{uncertainty}
    \end{itemize}
    }
\end{frame}

\begin{frame}{Who `ownes' or bears the Risk?}
    
    \mytext{
    \begin{itemize}
        \item Project Sponsor: risk \emph{of} the project
        \item Project Manager: risks \emph{in }the project
        \item Contractors?
        \item Government?
    \end{itemize}
    }
\end{frame}

\begin{frame}{Risk Treatment}
    
    \mytext{
    \begin{enumerate}
        \item \textbf{avoid}
        \item \textbf{reduce}
        \item \textbf{mitigate}
        \item \textbf{transfer}
        \item \textbf{share}
        \item \textbf{accept}
    \end{enumerate}
    }
\end{frame}

\begin{frame}{Risk Treatment per ISO 31000}
    
    \mytext{
    \begin{minipage}{0.65\textwidth}
        \begin{enumerate}
            \item \textbf{Avoiding} the risk by deciding not to start or continue with the activity that gives rise to the risk
            \item Accepting or increasing the risk in order to pursue an opportunity
            \item \textbf{Removing} the risk source
            \item \textbf{Changing} the \textbf{likelihood}
            \item \textbf{Changing} the \textbf{consequences}
            \item \textbf{Sharing} the risk with another party or parties (including contracts and risk financing)
            \item \textbf{Retaining} the risk by informed decision
        \end{enumerate}
    \end{minipage}
    }
\end{frame}

\begin{frame}{Hedging Risks}
    
    \mytext{
    \begin{itemize}
        \item \emph{Don't put all your eggs in one basket}
        \item Risk minimisation through diversification
        \item Have a `Plan B' (or multiple variants)
        \item Analyse core dependencies
        \item Install \textbf{buffers} to avoid \textbf{cascading}
        \item Insure, e.g.\ with liquidated damages
    \end{itemize}
    }
\end{frame}


\begin{frame}{Phases of Risk Treatment}
    
    \mytext{
\begin{enumerate}
    \item Risk Process initialisation
    \item Risk identification
    \item Qualitative risk assessment
    \item Quantitative risk assessment
    \item Risk response planning
    \item Risk communication
    \item Risk review
    \item Post-project review
\end{enumerate}
\vfill
    \cite[p.27]{hillson2012managing}
    }
\end{frame}

\begin{frame}{Risk Register}
    \showmytikz{../shared/tikz/tikz-PM-risk-register}
\end{frame}

% \input{break}



% \begin{frame}{A Bayesian View}
%     \begin{tikzpicture}[node distance=2cm]
%         \node[draw,rectangle] (past) at (0,0) {Past};
%         \node[draw,rectangle,right=2cm of past] (present) {Present};
%         \node[draw,rectangle,right of=present] (future2) [right=0.5cm]{Future 2};
%         \node[draw,rectangle,above of=future2] (future1) {Future 1};
%         \node[draw,rectangle,below of=future2] (future3) {Future 3};
%         \draw[->] (past) -- (present);
%         \draw[->] (present) -- node[above left]{25\%} (future1);
%         \draw[->] (present) -- node[above]{60\%}(future2);
%         \draw[->] (present) -- node[below left]{15\%}(future3);
%     \end{tikzpicture}
% \end{frame}




% \begin{frame}{Future Scenarios, Path Dependency}
%         \begin{tikzpicture}[yscale=0.9]
    
%         \node [shading = axis,rectangle, left color=red, anchor=south west, right color=green,shading angle=-45, minimum width=11cm, minimum height=6cm, opacity=0.4] (box) at (-1,-6){};
    
%         \node (ti) at (0,1) {t};
%         \node[draw,circle] (a) at (0,0) {A};
        
%         \node (tii) at (2,1) {t+1};
%         \node[draw,circle] (b) at (2,0) {B};
%         \node[draw,circle] (c) at (2,-1.2) {C};
        
%         \node (tiii) at (4,1) {t+2};
%         \node[draw,circle] (d) at (4,0) {D};
%         \node[draw,circle] (e) at (4,-1.7) {E};
%         \node[draw,circle] (f) at (4,-3.5) {F};
        
%         \node (tiv) at (6,1) {t+3};
%         \node[draw,circle] (g) at (6,0) {G};
%         \node[draw,circle] (h) at (6,-1.5) {H};
%         \node[draw,circle] (i) at (6,-3) {I};
%         \node[draw,circle] (k) at (6,-4.5) {K};
        
%         \node (tv) at (8,1) {t+4};
%         \node[draw,circle] (l) at (8,0) {L};
%         \node[draw,circle] (m) at (8,-1) {M};
%         \node[draw,circle] (n) at (8,-2.5) {N};
%         \node[draw,circle] (o) at (8,-4) {O};
%         \node[draw,circle] (p) at (8,-5) {P};        
%         \draw (a) -- (b);
%         \draw (a) -- (c);
%         \draw (b) -- (d);
%         \draw (b) -- (e);
%         \draw (c) -- (d);
%         \draw (c) -- (e);
%         \draw (c) -- (f);
%         \draw (d) -- (g);
%         \draw (d) -- (h);
%         \draw (e) -- (h);
%         \draw (e) -- (i);
%         \draw (f) -- (i);
%         \draw (f) -- (k);
%         \draw (g) -- (l);
%         \draw (g) -- (n);
%         \draw (h) -- (m);
%         \draw (h) -- (n);
%         \draw (i) -- (n);
%         \draw (i) -- (o);
%         \draw (k) -- (o);
%         \draw (k) -- (p);
%         \node (con) at (9,1) {6$^{\circ}$C};
%         \node (pro) at (9,-5) {1$^{\circ}$C};
%         \draw[thick,<->,>=latex] (con) -- (pro);
%    \end{tikzpicture}
% \end{frame}

% \begin{frame}{Inertia, Time Lags}
%         \begin{tikzpicture}[yscale=0.9]
    
%         \node [shading = axis,rectangle, left color=red, anchor=south west, right color=green,shading angle=-45, minimum width=11cm, minimum height=6cm, opacity=0.4] (box) at (-1,-6){};
    
%         \node (ti) at (0,1) {t};
%         \node[draw,circle] (a) at (0,0) {A};
        
%         \node (tii) at (2,1) {t+1};
%         \node[draw,circle] (b) at (2,0) {B};
%         \node[draw,circle] (c) at (2,-1.2) {C};
        
%         \node (tiii) at (4,1) {t+2};
%         \node[draw,circle] (d) at (4,0) {D};
%         \node[draw,circle] (e) at (4,-1.7) {E};
%         \node[draw,circle] (f) at (4,-3.5) {F};
        
%         \node (tiv) at (6,1) {t+3};
%         \node[draw,circle] (g) at (6,0) {G};
%         \node[draw,circle] (h) at (6,-1.5) {H};
%         \node[draw,circle] (i) at (6,-3) {I};
%         \node[draw,circle] (k) at (6,-4.5) {K};
        
%         \node (tv) at (8,1) {t+4};
%         \node[draw,circle] (l) at (8,0) {L};
%         \node[draw,circle] (m) at (8,-1) {M};
%         \node[draw,circle] (n) at (8,-2.5) {N};
%         \node[draw,circle] (o) at (8,-4) {O};
%         \node[draw,circle] (p) at (8,-5) {P};        
        
        
%         \draw (a) -- (b);
        
%             \draw[dashed,bend right=25] (a) to (d);
%             \draw[dashed,bend right=35] (a) to            (e);
%             \draw[dashed,bend right=35] (a) to (f);
        
%         \draw (a) -- (c);
        
%         \draw (b) -- (d);
%         \draw (b) -- (e);
%             \draw[dashed,bend left=25] (b) to (g);
%             \draw[dashed,bend right=5] (b) to (h);
        
%         \draw (c) -- (d);
%         \draw (c) -- (e);
%         \draw (c) -- (f);
        
%         \draw (d) -- (g);
%         \draw (d) -- (h);
        
%             \draw[dashed,bend left=25] (d) to (l);
%             \draw[dashed,bend right=5] (d) to (m);
        
%         \draw (e) -- (h);
%         \draw (e) -- (i);
        
%            \draw[dashed,bend left=15] (e) to (m);
%             \draw[dashed,bend right=5] (e) to (n);
%             \draw[dashed,bend right=25] (e) to (o);        
%         \draw (f) -- (i);
%         \draw (f) -- (k);
        
%             \draw[dashed,bend left=15] (f) to (n);
%             \draw[dashed,bend right=5] (f) to (o);
%             \draw[dashed,bend right=25] (f) to (p);        
        
%         \draw (g) -- (l);
%         \draw (g) -- (n);
        
%         \draw (h) -- (m);
%         \draw (h) -- (n);
        
%         \draw (i) -- (n);
%         \draw (i) -- (o);
        
%         \draw (k) -- (o);
%         \draw (k) -- (p);

%         \node (con) at (9,1) {6$^{\circ}$C};
%         \node (pro) at (9,-5) {1$^{\circ}$C};
        
%         \draw[thick,<->,>=latex] (con) -- (pro);

%    \end{tikzpicture}
% \end{frame}

\begin{frame}{Optimism Bias and its Correction}
    
    \mytext{
    \begin{minipage}{0.7\textwidth}
        \begin{description}
            \item[Optimism Bias (OB)] is the proven systematic tendency of people to be overly optimistic about the outcome of planned actions.
        \end{description}
        \begin{itemize}
            \item The ideal, most optimistic scenario is an extreme and probably  \emph{not realistic}
            \item Not just the project manager, but also contractors are prone to OB
            \item $\rightarrow$ Escalation!
            \item Type of planning error / planning risk
        \end{itemize}
    \end{minipage}
    }
\end{frame}

\begin{frame}{Optimism Bias and its Correction}
    
    \mytext{
    \begin{itemize}
        \item Planning errors due to OB likely to appear \emph{consistently!}
        \item Measure average OB, analyse distribution
        \item Review plan, correcting values!
        \item Introduce \textbf{buffers,} insurance, liquidated damages
        \item PM's buffer is different to contractor's buffer 
        \item[] $\rightarrow$ information bias is important
    \end{itemize}
    }
\end{frame}





% \begin{frame}{Be Not Afraid of Numbers and Statistics}
%     \begin{itemize}
%         \item Understand Key Performance Indicators  (KPIs)
%         \item Understand law of large numbers (LLN)
%         \item Understand Law of truly large numbers (LTLN)
%         \item Understand statistical significance
%         \item Understand Benford's law
%         \item Undstand false positives and false negatives
%         \item Understand Bayes' Law
%     \end{itemize}
% \end{frame}



% \begin{frame}{Understand Bayes' Law}
%     \begin{minipage}{0.5\textwidth}
%         \begin{equation}
%             P (A | B) = \frac{P(B|A) \cdot P (A)}{P(B)}
%         \end{equation}
    
%         \medskip
        
%         Typical example: testing
        
%         \medskip
        
%         \begin{description}
%             \item[$P(A)$] probability to have Covid-19
%             \item[$P(B)$] probability that test on Covid-19 is positive
%             \item[$P(A|B)$] probability to have Covid-19 \emph{given test is positive}
%             \item[$P(B|A)$] probability that test is positive \emph{given Covid-19}
%         \end{description}
%     \end{minipage}
% \end{frame}


% \begin{frame}{Uncertainty, Meta-view}
%     \begin{tikzpicture}[scale=0.5]
%         \draw (0,0.4) rectangle (10,10);
%         \node[anchor=west] at (10.5,4.8) {$P(\neg X)$};
%         \fill[draw, fill=red] (0,0.4) rectangle (10,0);
%         \node[anchor=west] at (10.5,0.2) {$P(X)$} ;
%     \end{tikzpicture}
    
%     A theory $A$ assigns the risk at $P(X)$
    
%     \framebreak
    
%     \begin{tikzpicture}[scale=0.5]
%         \draw (0,0.4) rectangle (10,10);
%         \node[anchor=west] at (10.5,4.8) {$P\left(\neg X \mid A\right)$};
%         \fill[draw, fill=red] (0,0.4) rectangle (10,0);
%         \node[anchor=west] at (10.5,0.2) {$P\left(X \mid A\right)$} ;  
%         \fill[draw, fill=green] (0,0) rectangle (3,10);
%         \node[anchor=south] at (1.5,10.5) {$P(\neg A)$};
%         \node[anchor=south] at (6.5,10.5) {$P(A)$};
%     \end{tikzpicture}
    
    
%     The \textbf{risk that the theory is wrong} is $P(\neg A)$, so the actual risk may be within the green area.
%     \cite{bostrom1977heinen}
% \end{frame}

% \input{break}

\begin{frame}{The Three Constraints}
    \showmytikz[caption={The Three Constraints (Risk)},label=fig:threeconstraints-risk]{../shared/tikz/tikz-PM-three-constraints-risk}
    \mycomment{
    

    Risk accompany costs, performance, time, relationships
    
    }
\end{frame}

\begin{frame}{Risk, Uncertainty (examples)}
    
    \mytext{
    \begin{description}
        \item[Time] 
        \begin{itemize}
            \item `student syndrom', procrastination
            \item Hofstadter's Law: It always takes longer than you expect, even when you take into account Hofstadter's Law
            \item \textbf{escalating} delays
            \item unforseen dependencies
            \item unforseen complications with other parties
            \item affects phases, milestones...
        \end{itemize}\end{description}
    }
\end{frame}

\begin{frame}{Risk, Uncertainty (examples)}
    
    \mytext{
\begin{description}
        \item[Budget] 
        \begin{itemize}
            \item hidden costs
            \item procurement risk
            \item scope creep
            \item delays
            \item interest rates
            \item financial / exchange rates etc.
            \item KPI, ROI etc. affected
        \end{itemize}\end{description}
    }
\end{frame}

\begin{frame}{Risk, Uncertainty (examples)}
    
    \mytext{
\begin{description}
        \item[Quality]
        \begin{itemize}
            \item supply problems
            \item uncertain requirements
            \item unclear assumptions
            \item environmental surprises
            \item performance
            \item security
            \item $\dots$
        \end{itemize}\end{description}
    }
\end{frame}

\begin{frame}{Risk, Uncertainty (examples)}
    
    \mytext{
\begin{description}
        \item[Relations]
        \begin{itemize}
            \item team change 
            \item project sponsor change
            \item customer change
            \item project manager personal situation
            \item contractor risk
            \item stakeholder risk
        \end{itemize}
    \end{description}
    }
\end{frame}


\begin{frame}{Realistic Promises}
    
    \mytext{
    \begin{itemize}
        \item `We will do this until tomorrow'
        \item `We will do this as fast as possible'
        \item Assign \textbf{resources first}
        \item Assign variance / bands
        \item `I will have two people working 8hs a day on this'
    \end{itemize}
    \medskip
    \begin{itemize}
        \item Fixed resources, variable outcome
        \item Fixed outcome... variable outcome.
        \item Set clear goals - set incentives
    \end{itemize}
    }
\end{frame}


\begin{frame}{Delays and Costs}
    
    \mytext{
    \begin{itemize}
        \item Delays mean higher costs
        \item Delays lead to delays
        \item Buying time almost impossible
        \item \textbf{Overlooked issues most dangerous}
    \end{itemize}
    }
\end{frame}



\begin{frame}{What is the Likely Outcome?}
    
    \mytext{
    \begin{itemize}
        \item Optimistic scenario probably \emph{not realistic}
        \item Importance to \textbf{know risk}
        \item Known probability distribution?
        \item Importance of \textbf{buffers,} insurance, liquidated damages
        \item Information biases: don't tell contractors about the buffer
    \end{itemize}
    }
\end{frame}




\begin{frame}{Risk and Failure}
    
    \mytext{
    \begin{itemize}
        \item Concept of sunk costs:
        \sloppypar{\emph{dont' throw good money after bad money}}
        \item But how to `pull out'? 
        \begin{itemize}
            \item obligations
            \item penalties
            \item reputation
            \item alternatives
        \end{itemize}
        \item There always is an alternative
        \item Scott's Antarctica mission, Mungo Park, Ford Edsel
        \item Importance of scenario planning in dynamic environments
        \item There \emph{are }alternatives
        \item Contingency plan(s)
    \end{itemize}
    }
\end{frame}

\begin{frame}{Roles and Responsibility}
    
    \mytext{
    \begin{itemize}
        \item Is the responsibility clear?
        \item Negative aspect of responsibility:
        \begin{itemize}
            \item In charge of doing somthing
            \item Responsible for consequences of failure
            \item Incentive to do something?
            \item Disincentive for failure?
        \end{itemize}
        \item Who bears what risk?
        \item Insurance? Moral hazard?
        \item Expertise? Consultancy?
        \item Shooting the messenger?
    \end{itemize}    
    }
\end{frame}

\begin{frame}{Moral Hazard}
    
    \mytext{
    \begin{itemize}
        \item Only the risk prone seek insurance
        \item Once insured, you take the risk
    \end{itemize}
    }
\end{frame}

\begin{frame}{Cascading Damages}
    
    \mytext{
    \begin{itemize}
        \item Actual damage outside your responsibility
        \item Tort?
        \item Liquidated damages?
        \item Redress? Compensation? Indemnification?
        \item Continguency plan
    \end{itemize}
    }
\end{frame}

\begin{frame}{Fault and Shame}
    
    \mytext{
    \begin{itemize}
        \item Concept of fault (and blame):
            \begin{itemize}
                \item internal attribution: a person's fault
                \item external attribution: circumstance
            \end{itemize}    
        \item Project manager to blame?
        \item \textbf{Shame free} environment
        \item Rewards for \textbf{solutions} to problems
        \item But: \textbf{incentives}
        \item Insurance, continguency plans
    \end{itemize}
    }
\end{frame}

\begin{frame}{Unknown Risk}
    
    \mytext{
    \begin{itemize}
        \item Worst risk is risk overlooked
        \item Allowances / reserves for the unknown
        \item Optimism bias
        \item Personal risk
        \item Own risk
        \item Partner failure
        \item Cascading problems
    \end{itemize}
    }
\end{frame}


\begin{frame}{Sum Up: For Your Project}
    
    \mytext{
    \begin{itemize}
        \item Show you are aware of risks
        \item Identify knowledge gaps
        \item Plan contingency
        \item List risks: risk register
        \item Quantify 
        \item Controlling
        \item Plan alternatives
        \item Plan buffers and insurances
        \item Plan project plan review cycles
    \end{itemize}
    }
\end{frame}

% \input{break}


% \begin{frame}{Transfer, Lasting Impact, Sustainability}
    
    
%     \begin{itemize}
%         \item How do you end a project?
%         \item How do you evaluate success?
%         \begin{itemize}
%             \item Good project management?
%             \item Good project?
%         \end{itemize}
%         \item How does the transition / handover happen?
%         \item What will be the lasting impact?
%         \item What risks do you pass on?
%         \item Is your project sustainable?
%     \end{itemize}
    
%     \framebreak
    
%         \resizebox{!}{0.8\textheight!}{
%             \input{tikz/tikz-port}
%         }

%     \framebreak
    
%     \begin{itemize}
%         \item plan closure!
%         \item document your work
%         \item gain approval
%         \item transfer ownership
%         \begin{itemize}
%             \item of profits
%             \item of dept and risks
%         \end{itemize}
%         \item dismiss employees
%         \item dissolve SPV
%         \item attention: lasting liabilities
%     \end{itemize}
% \end{frame}

% \begin{frame}{Sustainability and the Three Constraints}

%     \begin{tikzpicture}[scale=0.8]
%         \coordinate (A) at (90:2.5);
%         \coordinate (B) at (210:2.5);
%         \coordinate (C) at (330:2.5);
%         \coordinate (a) at (0:5);
%         \coordinate (b) at (120:5);
%         \coordinate (c) at (240:5);
    
        
%         \coordinate (X) at (0,-3.35);
%         \draw (A) -- (B) node[midway,below,rotate=60] {Costs};
%         \draw (B) -- (C) node[midway,above] {Quality};
%         \draw (C) -- (A) node[midway,below,rotate=-60] {Time};

%         % \draw (a) -- (b) node[midway,above,rotate=-30] {Economically Viable};
%         % \draw (b) -- (c) node[midway,above,rotate=90] {Socially Responsible};
%         % \draw (c) -- (a) node[midway,below,rotate=30] {Environmentally Sound};

        
%         \path[draw=blue, -stealth, rotate=90,
%       postaction={decorate,decoration={text effects along path,
%         text={\ Within good customer relations \ }, text align=center, 
%         text effects/.cd, 
%           text along path, 
%           every character/.style={fill=white, yshift=-0.5ex}}}]
%       (0:0) circle (4);
%     \end{tikzpicture}

%     \framebreak

%     \begin{tikzpicture}[scale=0.8]
%         \coordinate (A) at (90:2.5);
%         \coordinate (B) at (210:2.5);
%         \coordinate (C) at (330:2.5);
%         \coordinate (a) at (0:5);
%         \coordinate (b) at (120:5);
%         \coordinate (c) at (240:5);
    
        
%         % \coordinate (X) at (0,-3.35);
%         \draw (A) -- (B) node[midway,below,rotate=60] {Costs};
%         \draw (B) -- (C) node[midway,above] {Quality};
%         \draw (C) -- (A) node[midway,below,rotate=-60] {Time};

%         \draw (a) -- (b) node[midway,above,rotate=-30] {Economically Viable};
%         \draw (b) -- (c) node[midway,above,rotate=90] {Socially Responsible};
%         \draw (c) -- (a) node[midway,below,rotate=30] {Environmentally Sound};

        
%         \path[draw=blue, -stealth, rotate=90,
%       postaction={decorate,decoration={text effects along path,
%         text={\ Within good customer relations \ }, text align=center, 
%         text effects/.cd, 
%           text along path, 
%           every character/.style={fill=white, yshift=-0.5ex}}}]
%       (0:0) circle (4);
%     \end{tikzpicture}


% \end{frame}

% \begin{frame}{Sustainability}
%     \begin{itemize}
%         \item Strong sustainability
%         \item[] Act so that things remain the same
%         \item Weak sustainability
%         \item[] Substitute what you take by something of equal value for the future
%     \end{itemize}

%     \begin{quote}
%         "Sustainable development is development that meets the needs of the present without compromising the ability of future generations to meet their own needs"

%         \flushright Brundtland Report 1987
%     \end{quote}

%     \framebreak

%     \begin{description}
%         \item [Human Rights:] discrimination of vulnerable groups, civil rights, and fundamental rights and principles at work
%         \item [Labor Practices:] conditions of work, health and safety, and development and training
%         \item [The Environment:] sustainable resource use, pollution prevention, and climate change mitigation
%         \item [Fair Operating Practices:] anti-corruption, fair competition, and respect for property rights
%         \item [Consumer Issues:] fair contractual practices, dispute resolution, and fair marketing
%         \item [Community Involvement and Engagement:] employee training and skills development, wealth and income creation, and community involvement
%     \end{description}

% \cite{tharp2012}
    
% \end{frame}



% \begin{frame}{Project Closure Review}
%     \begin{itemize}
%         \item Lessons learned?
%         \item Satisfaction stakeholders?
%         \item Final balance?
%         \item Technical dept?
%         \item Heritage?
%         \item Follow-up?
%     \end{itemize}    
% \end{frame}

% \begin{frame}{Closure}
%     Project closure inclucedes...
%     \begin{itemize}
%         \item a number of milestones
%         \item a number of events
%         \begin{itemize}
%             \item public events
%             \item in-house / legal
%         \end{itemize}
%     \end{itemize}

%     \medskip

%     Project Closure requires time and attention!

%     \framebreak

%     Recall project phases:

%     \begin{itemize}
%         \item Pre-project phase
%         \item Initialisation phase
%         \item Actualisation phase
%         \item Closure phase
%         \item Post-closure phase?
%     \end{itemize}

%     Workload (for management) is unevenly distributed
    
%     Danger of crunch/burnout towards end of project!

%     Different specialists needed!
    
% \end{frame}

% \begin{frame}{Project Portfolio Management}
%     \begin{itemize}
%         \item Multiple projects at any time
%         \item PMO
%         \item Life Cycles of Projects complementary
%         \item Risk minimisation or accumulation?
%         \item Same PM/controlling software used?
%         \item \textbf{Choice} of projects: strategic goals
%         \item PPM is NOT a project
%         \item Introducing PPM IS a project
%     \end{itemize}
% \end{frame}


% \begingroup
%     \setbeamercolor{background canvas}{bg=violet!20}


%     \begin{frame}{Your Presentations}
%         \begin{itemize}
%             \item \textbf{Feb 15th 2025 9am }
%             \item ONLINE
%             \item State your topic on Aulis
%             \item 10 minutes suffice!
%         \end{itemize}
%     \end{frame}

%     \begin{frame}{Extra Date}
%         I am available on \textbf{Jan 28th for consultancy}
%     \end{frame}

%     \setbeamercolor{background canvas}{bg=green!20}

%     \begin{frame}{Module Closure}
%         \begin{itemize}
%             \item Thank you for your attention!
%             \item Wish you all the best!
%             \item Any questions - come forward!
%             \item Keep in touch!
%         \end{itemize}
%     \end{frame}

% \endgroup



% \begin{frame}{Project Governance Report}
%     \begin{itemize}
%         \item user requirements ('need' / reason / context)
%         \item solution design: the actual product
%         \item stakeholders, management team, identity, communication strategy
%         \item work breakdown structure (WBS)
%         \item work packages and responsibilities
%         \item (project plan: network / BPMN plan)
%         \item project plan: temporal plan / Gantt diagram
%         \item governance plan: incorporation, technique, controls  
%         \item documentation, controls and reporting
%         \item risk register and risk management
%         \item transfer, lasting impact, sustainability
%         \item requirements / business plan (overview)
%     \end{itemize}
% \end{frame}

